\documentclass[11pt]{article}
\usepackage{amsmath,amssymb,amsthm,mathtools}
\usepackage{algorithm}
\usepackage{algpseudocode}
\usepackage{geometry}
\geometry{margin=1in}
\usepackage{enumitem}
\setlist[enumerate]{leftmargin=*,label=(\alph*)}
\setlist[itemize]{leftmargin=*}

\newtheorem{theorem}{Theorem}
\newtheorem{lemma}{Lemma}
\newtheorem{corollary}{Corollary}

\newcommand{\prox}{\operatorname{prox}}
\newcommand{\R}{\mathbb{R}}
\newcommand{\inner}[2]{\left\langle #1,\, #2 \right\rangle}
\newcommand{\norm}[1]{\left\lVert #1 \right\rVert}

\title{Proximal Gradient Method: Algorithm, Convergence Theory, and Complete Proof}
\author{}
\date{}

\begin{document}
\maketitle

\section*{Algorithm Name}
Proximal Gradient Method (Forward--Backward Splitting) for minimizing a composite convex objective
\[
F(x) \coloneqq f(x) + g(x),
\]
where $f$ is convex and smooth, and $g$ is convex (possibly nonsmooth).

\section*{Mathematical Formulation}
Let $f:\R^d\to\R$ be convex and differentiable with $L$-Lipschitz gradient, and let $g:\R^d\to(-\infty,+\infty]$ be proper, closed, and convex. For a stepsize $\eta>0$, define the proximal operator
\[
\prox_{\eta g}(u) \;\coloneqq\; \arg\min_{x\in\R^d}\left\{ g(x) + \frac{1}{2\eta}\norm{x-u}^2 \right\}.
\]
The proximal gradient iteration is
\[
x_{k+1} \;=\; T_\eta(x_k) \;\coloneqq\; \prox_{\eta g}\bigl(x_k - \eta \nabla f(x_k)\bigr),
\]
i.e.,
\[
x_{k+1} \;=\; \arg\min_{x}\left\{ g(x) + \inner{\nabla f(x_k)}{x-x_k} + \frac{1}{2\eta}\norm{x-x_k}^2 \right\}.
\]
Hyperparameters:
- Stepsize $\eta \in (0, 1/L]$ for descent guarantees and rate.
- Alternative admissible range for nonexpansiveness: $\eta \in (0, 2/L)$.

\section*{Pseudocode}
\begin{algorithm}[H]
\caption{Proximal Gradient Method}
\begin{algorithmic}[1]
\State \textbf{Input:} Initial point $x_0\in\R^d$, stepsize $\eta>0$, maximum iterations $T$, tolerance $\varepsilon\ge 0$
\For{$k=0,1,2,\dots,T-1$}
  \State $g_k \gets \nabla f(x_k)$
  \State $y_k \gets x_k - \eta g_k$
  \State $x_{k+1} \gets \prox_{\eta g}(y_k)$
  \If{$\norm{x_{k+1}-x_k}\le \varepsilon$}
     \State \textbf{break}
  \EndIf
\EndFor
\State \textbf{Output:} $x_T$
\end{algorithmic}
\end{algorithm}

\section*{Dry Run Example}
We minimize $F(w)=f(w)+g(w)$ with $f(w)=(w-3)^2$ and $g\equiv 0$. Then $\nabla f(w)=2(w-3)$ and $\prox_{\eta g}=\mathrm{Id}$.
Choose $w_0=0$, $\eta=0.1$ and perform 3 iterations.
\[
w_{k+1} = w_k - \eta \nabla f(w_k) = w_k - 0.1\cdot 2(w_k-3) = 0.8\,w_k + 0.6.
\]
- Iteration 0: $w_0=0$, $g_0=2(0-3)=-6$, $w_1=0-0.1(-6)=0.6$, $f(w_1)=(0.6-3)^2=2.4^2=5.76$.
- Iteration 1: $w_1=0.6$, $g_1=2(0.6-3)=-4.8$, $w_2=0.6-0.1(-4.8)=1.08$, $f(w_2)=(1.08-3)^2=1.92^2=3.6864$.
- Iteration 2: $w_2=1.08$, $g_2=2(1.08-3)=-3.84$, $w_3=1.08-0.1(-3.84)=1.464$, $f(w_3)=(1.464-3)^2=1.536^2=2.359296$.

\section*{Assumptions}
Let $F(x)=f(x)+g(x)$ with:
- A1. $f:\R^d\to\R$ is convex and differentiable.
- A2. $\nabla f$ is $L$-Lipschitz: $\norm{\nabla f(x)-\nabla f(y)}\le L\norm{x-y}$ for all $x,y$.
- A3. $g:\R^d\to(-\infty,+\infty]$ is proper, closed, and convex.
- A4. The solution set $X^\star\coloneqq\operatorname{argmin}_x F(x)$ is nonempty; let $F^\star=\min_x F(x)$.
- A5. Stepsize $\eta\in(0,1/L]$ for the rate results; $\eta\in(0,2/L)$ suffices for nonexpansiveness and Fej\'er monotonicity.

\section*{Main Convergence Theorem}
\begin{theorem}[Global convergence and sublinear rate]\label{thm:main}
Under A1--A5, define the proximal gradient iterates $x_{k+1} = \prox_{\eta g}(x_k - \eta \nabla f(x_k))$.
Then for any $x^\star\in X^\star$:
\begin{enumerate}
\item (Fej\'er monotonicity via nonexpansiveness) For $\eta\in(0,2/L)$,
\[
\norm{x_{k+1}-x^\star}\le \norm{x_k - x^\star}\quad\text{for all }k\ge 0,
\]
and $(x_k)$ converges (weakly in Hilbert spaces, strongly in finite-dimensional spaces) to a point in $X^\star$.
\item (Objective descent) For $\eta\in(0,1/L]$,
\[
F(x_{k+1})\le F(x_k) - \left(\frac{1}{2\eta}-\frac{L}{2}\right)\norm{x_{k+1}-x_k}^2 \;\le\; F(x_k).
\]
\item (Rate with constants) For $\eta\in(0,1/L]$ and all $T\ge 1$,
\[
F(x_T)-F^\star \;\le\; \frac{\norm{x_0-x^\star}^2}{2\eta\,T} \;\le\; \frac{L\,\norm{x_0-x^\star}^2}{2\,T}.
\]
\end{enumerate}
\end{theorem}

\section*{Theoretical Properties}
- Stability and robustness: The iteration map $T_\eta=\prox_{\eta g}\circ (I-\eta\nabla f)$ is nonexpansive for $\eta\in(0,2/L)$, ensuring Fej\'er monotonicity of distances to the solution set.
- Hyperparameter sensitivity: 
  - If $\eta\in(0,1/L]$, $F(x_k)$ is nonincreasing and admits an $O(1/k)$ rate in function value.
  - If $\eta\in(1/L,2/L)$, nonexpansiveness holds but monotone decrease of $F$ need not hold.
  - If $\eta\ge 2/L$, $T_\eta$ may fail to be nonexpansive; instability may occur.
- Baselines:
  - If $g\equiv 0$, the method reduces to gradient descent with the same rate.
  - If $g$ is an indicator of a closed convex set, the method becomes projected gradient descent.

\section*{Discussion}
- Deterministic setting: There is no stochasticity; all expectations of random quantities reduce to their deterministic values. Formally, for any deterministic sequence $(x_k)$ and measurable $\phi$, $\mathbb{E}[\phi(x_k)]=\phi(x_k)$.

\section*{Preliminaries (Definitions and Lemmas)}
We work in $\R^d$ with the Euclidean norm.

\paragraph{Proximal operator and firm nonexpansiveness.}
\begin{lemma}[Firm nonexpansiveness of $\prox_{\eta g}$]\label{lem:firm}
For any proper closed convex $g$ and $\eta>0$, the proximal operator is firmly nonexpansive:
\[
\norm{\prox_{\eta g}(u)-\prox_{\eta g}(v)}^2 \;\le\; \inner{\prox_{\eta g}(u)-\prox_{\eta g}(v)}{u-v},\quad \forall u,v\in\R^d.
\]
Consequently, it is nonexpansive: $\norm{\prox_{\eta g}(u)-\prox_{\eta g}(v)}\le \norm{u-v}$.
\end{lemma}
\begin{proof}
Let $p=\prox_{\eta g}(u)$ and $q=\prox_{\eta g}(v)$. From the optimality condition of the proximal map,
\[
\frac{1}{\eta}(u-p)\in \partial g(p)\quad\text{and}\quad \frac{1}{\eta}(v-q)\in \partial g(q).
\]
By monotonicity of $\partial g$, $\inner{\frac{1}{\eta}(u-p)-\frac{1}{\eta}(v-q)}{p-q}\ge 0$. Multiplying by $\eta$ and expanding yields
\[
\inner{u-v}{p-q} \;\ge\; \norm{p-q}^2 \quad\text{[Step 1]},
\]
which is the firm nonexpansiveness inequality. The nonexpansiveness follows by Cauchy--Schwarz:
$\norm{p-q}^2 \le \norm{p-q}\,\norm{u-v} \Rightarrow \norm{p-q}\le \norm{u-v}$ \quad[Step 2].
\end{proof}

\paragraph{Smoothness, cocoercivity, and descent.}
\begin{lemma}[Descent lemma]\label{lem:descent}
If $\nabla f$ is $L$-Lipschitz, then for all $x,y\in\R^d$,
\[
f(y)\;\le\; f(x) + \inner{\nabla f(x)}{y-x} + \frac{L}{2}\norm{y-x}^2.
\]
\end{lemma}
\begin{proof}
Define $\phi(t) = f(x+t(y-x))$. Then $\phi'(t) = \inner{\nabla f(x+t(y-x))}{y-x}$ and $|\phi'(s)-\phi'(t)|\le L\norm{y-x}^2|s-t|$ by Lipschitzness. Integrating $\phi'$ from $0$ to $1$ and using the bound on $\phi'$ variations gives
\[
f(y)-f(x) = \int_0^1 \phi'(t)\,dt \le \phi'(0) + \int_0^1 L(1-t)\norm{y-x}^2\,dt = \inner{\nabla f(x)}{y-x} + \frac{L}{2}\norm{y-x}^2 \quad\text{[Step 3]}.
\]
\end{proof}

\begin{lemma}[Baillon--Haddad cocoercivity]\label{lem:bh}
If $f$ is convex and $\nabla f$ is $L$-Lipschitz, then for all $x,y\in\R^d$,
\[
\inner{\nabla f(x)-\nabla f(y)}{x-y} \;\ge\; \frac{1}{L}\,\norm{\nabla f(x)-\nabla f(y)}^2.
\]
\end{lemma}
\begin{proof}
We present a convex-conjugate proof. Since $\nabla f$ is $L$-Lipschitz, the convex conjugate $f^*$ is $(1/L)$-strongly convex (standard equivalence). Concretely, for all $p,q\in\R^d$,
\[
f^*(p) \;\ge\; f^*(q) + \inner{\nabla f^*(q)}{p-q} + \frac{1}{2L}\norm{p-q}^2 \quad\text{[Step 4]}.
\]
Set $p=\nabla f(x)$, $q=\nabla f(y)$. Using the identities $\nabla f^*(\nabla f(y))=y$ and $f^*(\nabla f(x)) = \inner{\nabla f(x)}{x} - f(x)$ (Fenchel--Young equality), we rewrite [Step 4] as
\[
\inner{\nabla f(x)}{x} - f(x) \;\ge\; \inner{\nabla f(y)}{y} - f(y) + \inner{y}{\nabla f(x)-\nabla f(y)} + \frac{1}{2L}\norm{\nabla f(x)-\nabla f(y)}^2 \quad\text{[Step 5]}.
\]
Rearrange terms:
\[
f(x) - f(y) \;\le\; \inner{\nabla f(x)}{x-y} - \inner{\nabla f(x)-\nabla f(y)}{y} - \frac{1}{2L}\norm{\nabla f(x)-\nabla f(y)}^2 \quad\text{[Step 6]}.
\]
Add and subtract $\inner{\nabla f(y)}{x}$ inside the RHS, then simplify to obtain
\[
f(x) - f(y) \;\le\; \inner{\nabla f(x)}{x} - \inner{\nabla f(y)}{y} - \inner{\nabla f(x)-\nabla f(y)}{y} - \frac{1}{2L}\norm{\nabla f(x)-\nabla f(y)}^2
= \inner{\nabla f(x)-\nabla f(y)}{x-y} - \frac{1}{2L}\norm{\nabla f(x)-\nabla f(y)}^2 \quad\text{[Step 7]}.
\]
By symmetry (swap $x,y$ in [Step 7]) and add the two inequalities to eliminate $f(x)-f(y)$, we get
\[
0 \;\le\; 2\,\inner{\nabla f(x)-\nabla f(y)}{x-y} - \frac{1}{L}\norm{\nabla f(x)-\nabla f(y)}^2 \quad\text{[Step 8]},
\]
which rearranges to the desired cocoercivity inequality.
\end{proof}

\begin{corollary}[Nonexpansiveness of gradient step]\label{cor:S}
Define $S_\eta(x)\coloneqq x - \eta \nabla f(x)$. For $\eta\in[0,2/L]$, $S_\eta$ is nonexpansive:
\[
\norm{S_\eta(x)-S_\eta(y)}\le \norm{x-y}\quad\forall x,y.
\]
\end{corollary}
\begin{proof}
Expand the square:
\begin{align*}
\norm{S_\eta(x)-S_\eta(y)}^2
&= \norm{x-y - \eta(\nabla f(x)-\nabla f(y))}^2 \\
&= \norm{x-y}^2 - 2\eta\,\inner{\nabla f(x)-\nabla f(y)}{x-y} + \eta^2 \norm{\nabla f(x)-\nabla f(y)}^2 \quad\text{[Step 9]}\\
&\le \norm{x-y}^2 - \left(\frac{2\eta}{L} - \eta^2\right)\norm{\nabla f(x)-\nabla f(y)}^2 \quad\text{(by Lemma \ref{lem:bh}) [Step 10]}\\
&\le \norm{x-y}^2 \quad\text{since } \frac{2\eta}{L}-\eta^2 \ge 0 \text{ for } \eta\in[0,2/L] \text{ [Step 11]}.
\end{align*}
Taking square roots yields the claim.
\end{proof}

\begin{corollary}[Nonexpansiveness of the iteration map]\label{cor:T}
Let $T_\eta = \prox_{\eta g}\circ S_\eta$. If $\eta\in[0,2/L]$, then $T_\eta$ is nonexpansive:
\[
\norm{T_\eta(x)-T_\eta(y)} \;\le\; \norm{x-y}\quad\forall x,y.
\]
\end{corollary}
\begin{proof}
By Lemma \ref{lem:firm}, $\prox_{\eta g}$ is nonexpansive; by Corollary \ref{cor:S}, $S_\eta$ is nonexpansive. Then
\[
\norm{T_\eta(x)-T_\eta(y)} = \norm{\prox_{\eta g}(S_\eta(x)) - \prox_{\eta g}(S_\eta(y))} \le \norm{S_\eta(x)-S_\eta(y)} \le \norm{x-y} \quad\text{[Step 12]}.
\]
\end{proof}

\begin{lemma}[Three-point inequality for the proximal subproblem]\label{lem:three-point}
Let $y\in\R^d$, $\eta>0$, and $x^+ \coloneqq \prox_{\eta g}(y)$. Then for any $x\in\R^d$,
\[
g(x^+) + \frac{1}{2\eta}\norm{x^+-y}^2 \;\le\; g(x) + \frac{1}{2\eta}\norm{x-y}^2 - \frac{1}{2\eta}\norm{x-x^+}^2.
\]
\end{lemma}
\begin{proof}
By optimality of $x^+$, for any $x$ and any subgradient $s^+\in\partial g(x^+)$,
\[
0 \in s^+ + \frac{1}{\eta}(x^+-y)\quad\Rightarrow\quad s^+ = \frac{1}{\eta}(y-x^+) \quad\text{[Step 13]}.
\]
Convexity of $g$ yields $g(x)\ge g(x^+) + \inner{s^+}{x-x^+}$ \quad[Step 14]. Substitute $s^+$ and complete the square:
\begin{align*}
g(x) &\ge g(x^+) + \frac{1}{\eta}\inner{y-x^+}{x-x^+} \\
&= g(x^+) + \frac{1}{2\eta}\bigl(\norm{x-y}^2 - \norm{x-x^+}^2 - \norm{x^+-y}^2\bigr) \quad\text{[Step 15]}.
\end{align*}
Rearranging gives the claim.
\end{proof}

\paragraph{Fixed-point characterization.}
\begin{lemma}[Optimality as a fixed point]\label{lem:fixed-point}
A point $x^\star$ is optimal for $F$ (i.e., $0\in \nabla f(x^\star) + \partial g(x^\star)$) if and only if
\[
x^\star \;=\; \prox_{\eta g}\bigl(x^\star - \eta \nabla f(x^\star)\bigr).
\]
\end{lemma}
\begin{proof}
Optimality condition $0\in \nabla f(x^\star) + \partial g(x^\star)$ is equivalent to $-\nabla f(x^\star)\in \partial g(x^\star)$. The proximal optimality for $u=x^\star - \eta \nabla f(x^\star)$ reads
\[
\frac{1}{\eta}(u-x^\star) = -\nabla f(x^\star) \in \partial g(x^\star) \quad\Rightarrow\quad x^\star=\prox_{\eta g}(u) \quad\text{[Step 16]}.
\]
The reverse implication follows by the same logic.
\end{proof}

\section*{Main Proof (Step-by-Step Derivation)}
We prove Theorem \ref{thm:main}.

\paragraph{Fej\'er monotonicity and convergence (Part 1).}
Let $x^\star\in X^\star$. By Lemma \ref{lem:fixed-point}, $x^\star$ is a fixed point of $T_\eta$, i.e., $x^\star=T_\eta(x^\star)$ \quad[Step 17]. Using Corollary \ref{cor:T},
\[
\norm{x_{k+1}-x^\star} = \norm{T_\eta(x_k)-T_\eta(x^\star)} \le \norm{x_k-x^\star} \quad\text{for } \eta\in(0,2/L) \quad\text{[Step 18]}.
\]
Thus $(\norm{x_k-x^\star})_{k\ge 0}$ is nonincreasing and bounded below by $0$, hence convergent. In $\R^d$, nonexpansive fixed-point iterations with a nonempty fixed-point set are Fej\'er monotone and every cluster point is a fixed point; thus the whole sequence converges to a point in $X^\star$ \quad[Step 19].

\paragraph{Objective descent (Part 2).}
Define $y_k \coloneqq x_k - \eta \nabla f(x_k)$. By Lemma \ref{lem:three-point} with $y=y_k$ and $x=x_k$,
\[
g(x_{k+1}) + \frac{1}{2\eta}\norm{x_{k+1}-y_k}^2 \;\le\; g(x_k) + \frac{1}{2\eta}\norm{x_k-y_k}^2 - \frac{1}{2\eta}\norm{x_k-x_{k+1}}^2 \quad\text{[Step 20]}.
\]
Note that $x_k-y_k=\eta \nabla f(x_k)$ and $x_{k+1}-y_k = x_{k+1}-x_k + \eta \nabla f(x_k)$ \quad[Step 21]. Hence
\[
g(x_{k+1}) \le g(x_k) + \inner{\nabla f(x_k)}{x_k - x_{k+1}} - \frac{1}{2\eta}\norm{x_{k+1}-x_k}^2 \quad\text{[Step 22]}.
\]
By Lemma \ref{lem:descent} with $x=x_k$ and $y=x_{k+1}$,
\[
f(x_{k+1}) \le f(x_k) + \inner{\nabla f(x_k)}{x_{k+1}-x_k} + \frac{L}{2}\norm{x_{k+1}-x_k}^2 \quad\text{[Step 23]}.
\]
Adding [Step 22] and [Step 23] yields
\begin{align*}
F(x_{k+1})
&\le F(x_k) - \left(\frac{1}{2\eta}-\frac{L}{2}\right)\norm{x_{k+1}-x_k}^2 \quad\text{[Step 24]}\\
&\le F(x_k) \quad\text{for }\eta\in(0,1/L] \quad\text{[Step 25]}.
\end{align*}

\paragraph{Rate (Part 3): $O(1/k)$ with constants.}
Apply Lemma \ref{lem:three-point} with $y=y_k$ and arbitrary $x\in\R^d$:
\[
g(x_{k+1}) + \frac{1}{2\eta}\norm{x_{k+1}-y_k}^2 \;\le\; g(x) + \frac{1}{2\eta}\norm{x-y_k}^2 - \frac{1}{2\eta}\norm{x-x_{k+1}}^2 \quad\text{[Step 26]}.
\]
Use Lemma \ref{lem:descent} again:
\[
f(x_{k+1}) \le f(x_k) + \inner{\nabla f(x_k)}{x_{k+1}-x_k} + \frac{L}{2}\norm{x_{k+1}-x_k}^2 \quad\text{[Step 27]}.
\]
Add [Step 26] and [Step 27], and insert $y_k=x_k-\eta\nabla f(x_k)$; then expand and simplify as in [Step 21]--[Step 24] to obtain
\begin{align*}
F(x_{k+1})
&\le f(x_k) + g(x) + \inner{\nabla f(x_k)}{x - x_k} + \frac{1}{2\eta}\norm{x - x_k}^2 - \frac{1}{2\eta}\norm{x - x_{k+1}}^2 + \left(\frac{L}{2}-\frac{1}{2\eta}\right)\norm{x_{k+1}-x_k}^2 \quad\text{[Step 28]}\\
&\le f(x_k) + \inner{\nabla f(x_k)}{x - x_k} + g(x) + \frac{1}{2\eta}\norm{x - x_k}^2 - \frac{1}{2\eta}\norm{x - x_{k+1}}^2 \quad\text{for }\eta\in(0,1/L] \text{ [Step 29]}\\
&\le f(x) + g(x) + \frac{1}{2\eta}\norm{x - x_k}^2 - \frac{1}{2\eta}\norm{x - x_{k+1}}^2 \quad\text{(convexity of $f$) [Step 30]}\\
&= F(x) + \frac{1}{2\eta}\norm{x - x_k}^2 - \frac{1}{2\eta}\norm{x - x_{k+1}}^2 \quad\text{[Step 31]}.
\end{align*}
Now choose $x=x^\star\in X^\star$ in [Step 31]:
\[
F(x_{k+1}) - F^\star \;\le\; \frac{1}{2\eta}\left(\norm{x^\star - x_k}^2 - \norm{x^\star - x_{k+1}}^2\right) \quad\text{[Step 32]}.
\]
Summing [Step 32] over $k=0$ to $T-1$ telescopes:
\begin{align*}
\sum_{k=0}^{T-1} \bigl(F(x_{k+1}) - F^\star\bigr)
&\le \frac{1}{2\eta}\left(\norm{x^\star - x_0}^2 - \norm{x^\star - x_T}^2\right)
\le \frac{1}{2\eta}\norm{x^\star - x_0}^2 \quad\text{[Step 33]}.
\end{align*}
By Part 2, $F(x_k)$ is nonincreasing for $\eta\in(0,1/L]$, hence $F(x_T)-F^\star \le \frac{1}{T}\sum_{k=1}^{T} (F(x_k)-F^\star)$ \quad[Step 34]. Combining with [Step 33],
\[
F(x_T)-F^\star \;\le\; \frac{1}{T}\cdot \frac{1}{2\eta}\norm{x^\star - x_0}^2 \;=\; \frac{\norm{x_0-x^\star}^2}{2\eta\,T} \quad\text{[Step 35]}.
\]
Finally, using $\eta\le 1/L$ gives $F(x_T)-F^\star \le \frac{L\norm{x_0-x^\star}^2}{2T}$ \quad[Step 36].

\paragraph{Expectation calculations.}
All quantities are deterministic; thus for any measurable $\psi$,
\[
\mathbb{E}[\psi(x_k)] = \psi(x_k) \quad\text{for all }k \quad\text{[Step 37]}.
\]
Therefore, the rate bounds hold both pointwise and in expectation.
\qed

\section*{Rate Derivation (With Constants)}
Collecting the constants from [Step 35]--[Step 36], under $\eta\in(0,1/L]$:
- Telescoping constant: $(2\eta)^{-1}$.
- Monotonicity yields last-iterate rate:
\[
F(x_T)-F^\star \le \frac{\norm{x_0-x^\star}^2}{2\eta\,T} \le \frac{L\,\norm{x_0-x^\star}^2}{2\,T}.
\]
- Descent per-iteration:
\[
F(x_k)-F(x_{k+1}) \ge \left(\frac{1}{2\eta}-\frac{L}{2}\right)\norm{x_{k+1}-x_k}^2.
\]

\section*{Special Cases}
- Smooth case $g\equiv 0$. Then $\prox_{\eta g}=\mathrm{Id}$ and $x_{k+1}=x_k-\eta\nabla f(x_k)$. The same proof (with $g\equiv 0$) gives
\[
f(x_T)-f^\star \le \frac{\norm{x_0-x^\star}^2}{2\eta\,T} \le \frac{L\norm{x_0-x^\star}^2}{2T},\quad \eta\in(0,1/L].
\]
- Projected gradient: If $g=\iota_C$ is the indicator of a nonempty closed convex set $C$, then $\prox_{\eta g}=\Pi_C$, the Euclidean projection; the method becomes $x_{k+1}=\Pi_C(x_k-\eta\nabla f(x_k))$ with the same rate.
- Strong convexity (informal). If $F$ is $\mu$-strongly convex and $\eta\in(0,1/L]$, a linear rate holds:
\[
F(x_k)-F^\star \le \left(1-\eta\mu\right)^k\bigl(F(x_0)-F^\star\bigr).
\]
A full derivation proceeds by strengthening [Step 32] using strong convexity to bound $F(x_{k+1})-F^\star$ in terms of $\norm{x_{k+1}-x^\star}^2$ and employing Fej\'er monotonicity.

\end{document}